\documentclass{article}

\textwidth=6.5in
\textheight=8.5in
\voffset=-54pt
\oddsidemargin=0in
\evensidemargin=0in
\setlength{\parskip}{8pt}
\setlength{\parindent}{0pt}

\usepackage{amsmath, amssymb}
\usepackage{ifthen}

\usepackage[inline]{enumitem}
\setlist{topsep=0pt, parsep=8pt}

\usepackage[rgb,table]{xcolor}\convertcolorsUfalse
\usepackage{graphicx}

% change hyperref options
\PassOptionsToPackage{hyphens}{url}
\usepackage[bookmarks,colorlinks,linkcolor=black,citecolor=blue,pdfstartview=FitH,urlcolor=blue]{hyperref}
\hypersetup{pdfpagemode=UseNone}
\usepackage{bookmark}

\usepackage{graphicx}
\usepackage{multicol}
\usepackage{framed}
\usepackage{array}

\usepackage{icomma}

\usepackage{pifont}

\usepackage{soul}
\usepackage[normalem]{ulem}

\usepackage{pgfplots}
\pgfplotsset{compat=1.18}  
\pgfplotscreateplotcyclelist{mycolor}{black, blue, red, green!70!black, magenta, purple, cyan, orange }  
\pgfplotsset{
  disabledatascaling,
  cycle list=tikzcolor, 
  axis lines=middle,
  every axis plot/.style={no marks, thick},
  every axis/.style={
    enlarge x limits=false,
    enlarge y limits=auto,
    xlabel style={at={(current axis.right of origin)},anchor=west},
    ylabel style={at={(current axis.above origin)},anchor=south},
    % extra x and y ticks for asymptotes
    extra x tick style={grid=major, grid style={dashed,black}},
    extra y tick style={grid=major, grid style={dashed,black}},
    /pgf/number format/1000 sep={},
  },    
  tick label style = {font=\footnotesize, fill=\currentbgcolor, inner sep=0pt},    
  xticklabel shift=1pt,    
  scaled ticks=false,
  scaled y ticks=false,
  unbounded coords=jump,
  samples=101,
  minor tick length=0,
  scatter/classes={
    open={mark=*,draw=black,fill=white, mark size=2, thin},
    closed={mark=*,draw=black,fill=black, mark size=2, thin}
  },
  width=3in,
}
\pgfplotsset{open/.style={only marks, mark=*, fill=white, thin, mark size=2}}
\pgfplotsset{closed/.style={only marks, mark=*, draw=black, fill=black,
    thin, mark size=2}}
\pgfplotsset{labeled/.style={nodes near coords, only marks, mark=*, draw=black, fill=black, thin, mark size=2, point meta=explicit symbolic}}

\pgfplotsset{gridgraph/.style={clip=false, axis line style={shorten
      >=-8pt}, xlabel style={xshift=8pt}, ylabel style={yshift=8pt},
    enlargelimits=false, grid=major, tick style={black}}} 

\newcommand{\axes}[1][]{
  \begin{tikzpicture}
    \begin{axis}[xtick=\empty, ytick=\empty, ymin=-2, ymax=2,
      xmin=-4.25, xmax=4.25, #1]
    \end{axis}
  \end{tikzpicture}
}
\usepackage{tikz}
\usetikzlibrary{
  matrix,%
  % enables the snake paths
  decorations.pathmorphing,%
  % enables bigger arrows
  decorations.markings,%
  decorations.pathreplacing,%
  decorations.text,%
  calc,%
  patterns,%
  shapes.geometric,%
  arrows,
  decorations.shapes,
  positioning,
  plotmarks,
  intersections
}
\tikzset{>=latex} % sets default arrow type in tikz
\tikzset{font=\normalsize}
\tikzset{mark size=1.5pt, mark options=thin}
\tikzset{pin distance=4pt,
  every pin edge/.style={<-, thin, shorten <= -2pt}}

\interfootnotelinepenalty=10000
\renewcommand{\thefootnote}{(\arabic{footnote})}

\usepackage{multirow, bigdelim}

% change to \usepackage[sols]{handout} to see solutions
\usepackage[sols]{handout}

\providecommand{\check}{}
\renewcommand{\check}{\ifmmode{\text{ \ding{51} }}\else{\ding{51} }\fi}

\newcommand\iscurrentenumi[1]{%
  \ifthenelse{\equal{#1}{\theenumi}}{}{#1}}
\setenumerate[1]{ref=\#\arabic*}
\setenumerate[2]{ref=\protect\iscurrentenumi{\theenumi}(\alph*)}
\newcommand\enumiiprefix[2]{%
  \ifthenelse{{\equal{#1}{\theenumi}}\AND{\equal{#2}{(\alph{enumii})}}}{}{}%
  \ifthenelse{{\equal{#1}{\theenumi}}\AND{\NOT{\equal{#2}{(\alph{enumii})}}}}{#2}{}%
  \ifthenelse{\equal{#1}{\theenumi}}{}{#1#2}%
}
\setenumerate[3]{ref=\protect\enumiiprefix{\theenumi}{(\alph{enumii})}\roman*}

\newcommand{\related}[1]{%
\fcolorbox{black}{yellow}{%
  %\parbox[t]{12pt}{$\clubsuit$}%
  \parbox[t]{\linewidth-8pt}{\emph{Related problems:} #1}}}

\renewcommand{\answer}[1]{#1}

\definecolor{purple}{rgb}{0.5,0,1.0}
\definecolor{hotpink}{rgb}{1, 0, 0.5}

\definecolor{skyblue}{rgb}{0, 0.5, 1}

\usepackage{amsthm}

% make URLs print in the same font
\usepackage{url}
\urlstyle{same}

\usepackage{pifont}
\def\exend{\hfill\ding{118}}

%\makeatletter\@addtoreset{equation}{section}\makeatother
%\renewcommand{\theequation}{\arabic{section}.\arabic{equation}}

\newtheoremstyle{theorem}% name
{8pt}%      Space above
{0pt}%      Space below
{\itshape}%         Body font
{}%         Indent amount (empty = no indent, \parindent = para indent)
{\bfseries}% Thm head font
{.}%        Punctuation after thm head
{.5em}%     Space after thm head: " " = normal interword space;
                                %       \newline = linebreak
{}%         Thm head spec (can be left empty, meaning `normal')

\theoremstyle{theorem}

\let\Theorem\undefined
\newtheorem{Theorem}[equation]{Theorem}
\let\Definition\undefined
\newtheorem{Definition}[equation]{Definition}
\let\Summary\undefined
\newtheorem{Summary}[equation]{Summary}
\let\Conjecture\undefined
\newtheorem{Conjecture}[equation]{Conjecture}
\let\Fact\undefined
\newtheorem{Fact}[equation]{Fact}

\renewcommand{\thesubsection}{\arabic{subsection}}
\makeatletter
\def\@seccntformat#1{\@ifundefined{#1@cntformat}%
   {\csname the#1\endcsname\quad}%       default
   {\csname #1@cntformat\endcsname}}%    enable individual control
\newcommand\section@cntformat{}
\makeatother

  \usepackage{exam}

\usepackage{fancyhdr}
\lhead{\textcolor{white}{This content is protected and may not be shared, uploaded, or distributed.}}
\rhead{}
\rfoot{\vspace*{\fill}\footnotesize\textcopyright{} \emph{Harvard Math Ma}}
\renewcommand{\headrulewidth}{0pt}
\pagestyle{fancy}
\begin{document}

  \title{Practice Mini-Exam 2}

\instructions{instructions.tex}{}

\listofproblems

\begin{enumerate}

\item \points{12} \textbf{Algebra.}

  \begin{enumerate}
  \item \label{fall2015-1a-practice-gateway-12} 

    \begin{problem}[\vfill]
      Rewrite $\displaystyle \frac{a-1}{a+1} + \frac{2a^2}{a^2-1}$ as a single simplified fraction.
    \end{problem}

    \begin{solution}
      A common denominator is $a^2 - 1 = (a - 1)(a + 1)$:
      \begin{align*}
        \frac{a-1}{a+1} + \frac{2a^2}{a^2-1} &= \frac{(a - 1)(a - 1)}{(a +
          1)(a - 1)} + \frac{2a^2}{a^2 - 1} \\
        &= \frac{a^2 - 2a + 1 + 2a^2}{a^2 - 1} \\
        &= \boxed{\frac{3a^2 - 2a + 1}{a^2 - 1}}
      \end{align*}
    \end{solution}

  \item \label{fall2017-practice-gateway-3-11} 

    \begin{problem}[\vfill]
      Solve $ y^2 - 4y = y^2(7 + 2y)$ for $y$.
    \end{problem}

    \begin{solution}
      We can rewrite this as 
      \begin{align*}
        y^2 - 4y & = 7y^2 + 2y^3 \\
        \intertext{This isn't linear in $y$  (if you aren't sure what this means, see \href{https://people.math.harvard.edu/~jjchen/mathma/handouts/solving-equations.html}{the ``Algebra Practice: Solving equations'' handout}), so we'll move everything to one side and try to factor:} %
        0 & = 2y^3 + 6y^2 + 4y \\
        0 & = 2y(y^2 + 3y + 2) \\
        0 & = 2 y(y + 1)(y + 2) 
      \end{align*}
      In order for a product of terms to be 0, one of the terms must be 0, so $y = \boxed{0, -1, -2}$.
    \end{solution}

  \item \label{fall2017-practice-gateway-3-10-modified} 

    \begin{problem}[\vfill\newpage]
      Simplify the following expression, assuming $c, y > 0$:
      \[
        (y - bc)(by + c) - \frac{\sqrt[3]{8c^3 y^{12}}}{(2y)^3}
      \]
    \end{problem}

    \begin{solution}
      Let's work on the two pieces separately first:
      \begin{align*}
        (y - bc)(by + c) & = by^2 -b^2cy + cy - bc^2\\
        \frac{\sqrt[3]{8c^3 y^{12}}}{(2y)^2} &= \frac{\left( 8 c^3 y^{12} \right)^{1/3}}{2^3 y^3} \\
        &= \frac{8^{1/3} c^{3/3} y^{12/3}}{8 y^3} \\
        &= \frac{2 c y^4}{8y^3} \\
        &= \frac{1}{4} c y \\
        \shortintertext{So,}
        (y - bc)(by + c) - \frac{\sqrt[3]{8c^3 y^{12}}}{(2y)^3} &= by^2 -b^2cy + cy - bc^2 - \frac{1}{4} c y \\
        &= \boxed{by^2 -b^2cy + \frac{3}{4} cy - bc^2}
      \end{align*}

    \end{solution}

  \item \label{fall2018-algebra-gateway-v2-3} 

    \begin{problem}[\vfill\newpage]
      Solve $\displaystyle \frac{2Y+1}{P} + \frac{Y}{P+3}=0$ for $P$.
    \end{problem}

    \begin{solution}
      Let's first get rid of the fractions by multiplying both sides of the equation by $P(P + 3)$:
      \begin{align*}
        P (P + 3) \frac{2Y+1}{P} + P (P + 3) \frac{Y}{P + 3} &= 0 \\
        (P + 3)(2Y + 1) + PY &= 0 \\
        \intertext{This is linear in $P$ (if you aren't sure what this means, see \href{https://people.math.harvard.edu/~jjchen/mathma/handouts/solving-equations.html}{the ``Algebra Practice: Solving equations'' handout}), so we'll try to rewrite it as $(\text{something}) P = \text{something}$.} %
        P (2Y + 1) + 3 (2Y + 1) + PY &= 0 \\
        (2Y + 1) P + YP &= - 3(2Y + 1) \\
        (3Y + 1) P &= -3(2Y + 1) \\
        P &= \boxed{-\frac{3 (2Y + 1)}{3Y + 1}}
      \end{align*}
    \end{solution}

  \end{enumerate}

\item \points{5}  % 2.1.3

  The graph of $f(t)$ is the line drawn below. Let $A(x)$ be the area between the line and the $t$-axis from $t = -2$ to $t = x$, where $x$ is in the interval $[-2, 8]$. Two examples are shown below.
  \begin{center}
    \begin{tikzpicture}[declare function={f(\x)=-(\x-8)/2.5;}] 
      \begin{axis}[axis equal image, xtick={-2, 2, 8}, xticklabels={$-2$, $x$, $8$}, ytick=\empty, xlabel=$t$, width=2.5in, clip=false]
        \addplot+[domain=-3:9] {f(x)} node[pos=0, left]{$y = f(t)$}; %
        \addplot+[fill=black!50!white, fill opacity=0.5, domain=-2:2, draw=none] {f(x)} \closedcycle; % 
      \end{axis}
    \end{tikzpicture}
    \hspace{0.25in}
    \begin{tikzpicture}[declare function={f(\x)=-(\x-8)/2.5;}] 
      \begin{axis}[axis equal image, xtick={-2, 5, 8}, xticklabels={$-2$, $x$, $8$}, ytick=\empty, xlabel=$t$, width=2.5in, clip=false]
        \addplot+[domain=-3:9] {f(x)} node[pos=0, left]{$y = f(t)$}; %
        \addplot+[fill=black!50!white, fill opacity=0.5, domain=-2:5, draw=none] {f(x)} \closedcycle; % 
      \end{axis}
    \end{tikzpicture}
  \end{center}

  \begin{enumerate}

  \item
    \begin{problem}[\vfill]
      Is $A(x)$ a function? Explain your reasoning.
    \end{problem}

    \begin{solution}
      For each $x$ in $[-2, 8]$, there is a single value for the area $A(x)$, so $A(x)$ \fbox{is a function}. 
    \end{solution}

  \item
    \begin{problem}[\vfill]
      If $A(x)$ is a function, when is it increasing? When is it decreasing? Explain briefly. 
    \end{problem}

    \begin{solution}
      $A(x)$ is \fbox{increasing on $[-2, 8]$}: as $x$ increases, the area we're looking at also increases. 
    \end{solution}

  \item
    \begin{problem}[\vfill\newpage]
      If $A(x)$ is a function, is it linear, concave up, or concave down? Explain briefly. 
    \end{problem}

    \solonly{\related{\#3, {{\#4}} on \href{https://people.math.harvard.edu/~jjchen/mathma/docs/worksheets/worksheet01-sols.html}{the worksheet ``Introduction to Functions''}; \href{https://people.math.harvard.edu/~jjchen/mathma/docs/homework/PS01.html}{Problem Set 1}, {{\#2}}}}

    \begin{solution}
      If we think of $x$ increasing by a small amount, say by 0.5, from $-2$, then $A(x)$ increases by the blue area below. But if $x$ increases again by that same small amount, then $A(x)$ increases by the pink area below, which is smaller than the blue area.
      \begin{center}
        \begin{tikzpicture}[declare function={f(\x)=-(\x-8)/2.5;}] 
          \begin{axis}[axis equal image, xtick={-2, 8}, ytick=\empty, xlabel=$t$, width=2.5in, clip=false]
            \addplot+[domain=-3:9] {f(x)} node[pos=0, left]{$y = f(t)$}; %
            \addplot+[fill=skyblue, fill opacity=0.5, domain=-2:-1.5, draw=none] {f(x)} \closedcycle; % 
            \addplot+[fill=hotpink, fill opacity=0.5, domain=-1.5:-1, draw=none] {f(x)} \closedcycle; % 
          \end{axis}
        \end{tikzpicture}
      \end{center}
      If we keep increasing $x$ by the same small amount, the additional area will keep shrinking. This means the graph of $A$ will have this shape:
      \begin{center}
        \begin{tikzpicture}
          \begin{axis}[ytick=\empty, height=2in, xtick={-2, 8}]
            \addplot+[domain=-2:8] { -0.2*x^2+3.2*x + 7.2}; 
          \end{axis}
        \end{tikzpicture}
      \end{center}
      So, the graph of $A$ is \fbox{concave down}.      
    \end{solution}
  \end{enumerate}

\item \points{6} 

  In each part, you're given a list of properties. Give an example of a function $f(x)$ that satisfies \uline{all} of the properties; please give a formula for $f$. 

  \emph{Hint:} Can you alter one of our familiar functions? 

  \begin{enumerate}

  \item
    \begin{problem}[\vfill]
      \begin{itemize}[nosep]
      \item $f$ has domain $(-\infty, 2) \cup (2, \infty)$
      \item $f$ has a vertical asymptote at $x = 2$
      \item $f$ has upper bound $-3$
      \end{itemize}
    \end{problem}

    \begin{solution}
      A function from our basic library that seems helpful is $\dfrac{1}{x^2}$ because it has a vertical asymptote and a lower bound of 0. If we shift it right by 2 units, flip it over the $x$-axis, and move it down by $3$, then it will satisfy all of the given properties. The formula for this is $\boxed{f(x) = - \dfrac{1}{(x - 2)^2} - 3}$. Here's the graph of this example:
      \begin{center}
        \begin{tikzpicture}
          \begin{axis}[xtick={2}, ymax=1, ytick={-3}, enlarge y limits=false]
            \addplot+[domain=-1:5, restrict y to domain=-12:-3] {-1/(x-2)^2 - 3};
            \addplot+[domain=-1:5, dashed] {-3};
            \addplot+[domain=-10.75:1, dashed] (2, x); 
          \end{axis}
        \end{tikzpicture}
      \end{center}
    \end{solution}

  \item
    \begin{problem}[\vfill\newpage]
      \begin{itemize}[nosep]
      \item $f$ has domain $[2, \infty)$
      \item $f$ is continuous, decreasing without bound, and concave down on $[2, \infty)$.
      \end{itemize}
    \end{problem}

    \begin{solution}
      A function from our basic library that seems helpful is $\sqrt{x}$; if we shift it right by 2 units and flip it over the $x$-axis, that will have the behavior we want. The formula for this is $\boxed{f(x) = - \sqrt{x - 2}}$, and the graph of this example looks like this:
      \begin{center}
        \begin{tikzpicture}
          \begin{axis}[xmin=0, xtick={2}, ytick=\empty, ymax=2, enlarge y limits=false, width=2.5in]
            \addplot+[domain=2:6] {- sqrt(x - 2)};%
            \addplot+[closed] coordinates { (2, 0) }; %
          \end{axis}
        \end{tikzpicture}
      \end{center}
    \end{solution}

  \end{enumerate}

\item \label{fall2017-premidterm-7} \points{7}  % JC

  One afternoon, Jorge decides to go for a long run. Here is a graph of $f(t)$, the distance in miles Jorge has run by $t$ hours after noon. The graph is linear on $[1, 2]$ and $[2.5, 3]$.

  \begin{tikzpicture}
    \begin{axis}[xmin=0, xmax=3, ytick={2, 4, ..., 14}, gridgraph, grid=both, 
      minor
      xtick={0.5, 1, ..., 3}, minor ytick={1, 2, ..., 13}, ymax=14, xlabel=$t$, ylabel=$f(t)$, 
      ]
      \addplot+[only marks] coordinates {(1, 0) (3, 13)}; %
      \addplot+[sharp plot, black, very thick] coordinates {(1, 0) (2, 8)}; %
      \addplot+[domain=2:2.5, black, very thick] {-24 + 24*x - 4*x^2}; %
      \addplot+[sharp plot, black, very thick] coordinates {(2.5, 11) (3, 13)}; % 
    \end{axis}
  \end{tikzpicture}

  \begin{enumerate}
  \item
    \begin{problem}[\vfill]
      What time did Jorge's run start?
    \end{problem}

    \begin{solution}
      Jorge started at $t = 1$, which is 1 hour after noon. So, Jorge started running at \fbox{1 pm}.
    \end{solution}

  \item
    \begin{problem}[\vfill]
      What time did Jorge's run end?
    \end{problem}

    \begin{solution}
      \answer{3 pm}. 
    \end{solution}

  \item
    \begin{problem}[\vfill]
      What was Jorge's average speed over his entire run? Please include units in your answer.
    \end{problem}

    \begin{solution}
      $\text{average speed} = \dfrac{\text{distance traveled}}{\text{time taken}} = \dfrac{13 \text{ miles}}{2 \text{ hours}} = \boxed{6.5 \text{ mph}}$. 
    \end{solution}

  \item
    \begin{problem}[\vfill]
      What was Jorge's speed at 1:30 pm? Include units.
    \end{problem}

    \begin{solution}
      Since the graph is linear on $[1, 2]$, Jorge's speed was constant in that time. So, his speed at 1:30 pm is the same as his average speed between $t = 1$ and $t = 2$, which is $\dfrac{8 \text{ miles}}{1 \text{ hour}} = \boxed{8 \text{ mph}}$. 
    \end{solution}

  \item Write expressions for the following \textbf{in terms of $\bm{f}$}.

    \begin{enumerate}
    \item
      \begin{problem}[\vfill]
        The distance in miles Jorge has run $s$ minutes after noon.
      \end{problem}

      \begin{solution}
        $s$ minutes after noon is the same as $\dfrac{s}{60}$ hours after noon (for example, 120 minutes = 2 hours), so the distance Jorge has run by this time is $\boxed{f \left( \frac{s}{60} \right)}$. 
      \end{solution}

    \item
      \begin{problem}[\vfill]
        The distance in miles Jorge has run $h$ hours into his run.
      \end{problem}

      \begin{solution}
        Jorge's run starts at 1 pm, so $h$ hours into his run is $h + 1$ hours after noon. Therefore, the distance Jorge has run by this time is $\boxed{f(h + 1)}$.

        \check A good strategy to check your answer is to try a concrete example. For example, $h = 0.5$ is 1:30 pm, so his distance then is $f(1.5)$.
      \end{solution}

    \item
      \begin{problem}[\vfill\newpage]
        The distance in km Jorge has run $v$ hours after noon. (1 mile is 1.6 km.)
      \end{problem}

      \begin{solution}
        $v$ hours after noon, Jorge has run $f(v)$ miles, which is the same as $\boxed{1.6 f(v)}$ km. 
      \end{solution}
    \end{enumerate}

  \end{enumerate}

\item \label{fall2019-premidterm-practice2-2} \points{7}  \finishexam

  Below is the graph of a function $f(x)$.
  \begin{center}
    \begin{tikzpicture}[declare function={f(\x)=-.5*.125*\x*\x*\x*\x+.5*2.5*\x*\x-4;}]
      \begin{axis}[gridgraph, grid=both, minor xtick={-12, -11, ..., 12}, minor ytick={-12, -11, ..., 12}, xtick={-12, -10, ..., 12}, ytick={-12, -10, ..., 12}, xmin=-12, xmax=12, ymin=-12, ymax=12, axis equal image, width=5in]
        \begin{scope}[very thick]
          \addplot+[domain=-4.5:4.5]{f(x)};%
          \addplot+[sharp plot] coordinates { (-4.5, {f(-4.5)}) (-6, -6)};%
          \addplot+[sharp plot] coordinates { (4.5, {f(-4.5)}) (6, -6)};%
        \end{scope}
        \addplot+[closed] coordinates {(-6, -6) (6, -6)}; 
      \end{axis}
    \end{tikzpicture}
  \end{center} 

  \begin{enumerate}

  \item
    \begin{problem}[\vfill]
      Is $f(x)$ an even function, an odd function, or neither? How can you tell from the graph? 
    \end{problem}

    \begin{solution} 
      The graph of $f$ is symmetric about the $y$-axis, so $f$ is \fbox{even}.
    \end{solution}

  \item
    \begin{problem}[\vfill\newpage]
      Is $g(x) = f(x)+3$ an even function, an odd function, or neither? How do you know?
    \end{problem}

    \begin{solution}
      The function $f(x)+3$ is a vertical shift of the function $f(x)$ up by $3$ units. After the vertical shift, the graph is still symmetrical about the $y$-axis, so $g$ is \fbox{even}. 
    \end{solution} 

  \item
    \begin{problem}[\vfill]
      What's the domain of $f(x)$? 
    \end{problem}

    \begin{solution}
      \answer{$[-6, 6]$}
    \end{solution}

  \item
    \begin{problem}[\vfill]
      What's the domain of $h(x) = f(2x)$? 
    \end{problem}

    \begin{solution}
      We need $2x$ to be in the domain of $f$. This means that $-6 \leq 2x \leq 6$, so $-3 \leq x \leq 3$. So, the domain of $h$ is $\boxed{[-3, 3]}$. 
    \end{solution} 

  \item
    \begin{problem}
      Sketch a graph of $h(x) = f(2x)$.
    \end{problem}

    \probonly{\axes[gridgraph, grid=both, minor xtick={-12, -11, ..., 12}, minor ytick={-12, -11, ..., 12}, xtick={-12, -10, ..., 12}, ytick={-12, -10, ..., 12}, xmin=-12, xmax=12, ymin=-12, ymax=12, axis equal image, width=5in]}

    \begin{solution}
      \begin{center}
        \answer{
          \begin{tikzpicture}[declare function={f(\x)=-.5*.125*\x*\x*\x*\x+.5*2.5*\x*\x-4;}]
            \begin{axis}[gridgraph, grid=both, minor xtick={-12, -11, ..., 12}, minor ytick={-12, -11, ..., 12}, xtick={-12, -10, ..., 12}, ytick={-12, -10, ..., 12}, xmin=-12, xmax=12, ymin=-12, ymax=12, axis equal image, width=5in]
              \begin{scope}[very thick]
                \addplot+[domain=-2.25:2.25]{f(2*x)};%
                \addplot+[sharp plot] coordinates { (-2.25, {f(-4.5)}) (-3, -6)};%
                \addplot+[sharp plot] coordinates { (2.25, {f(-4.5)}) (3, -6)};%
              \end{scope}
              \addplot+[closed] coordinates {(-3, -6) (3, -6)}; 
            \end{axis}
          \end{tikzpicture}
        }
      \end{center}
    \end{solution} 

  \item
    \begin{problem}[\vfill\newpage]
      Is $j(x) = f(x+2)$ an even function, an odd function, or neither? How do you know?
    \end{problem}

    \begin{solution} 
      The graph of $j$ is a horizontal shift of the graph of $f(x)$ 2 units to the left. This translation will ruin the symmetry about the $y$-axis, so $j$ is \fbox{neither even nor odd}.
    \end{solution} 
  \end{enumerate} 

\end{enumerate}

\end{document}
